\subsubsection{Traces of minimizing metrics}
\label{sec:warped-minimizer-trace-limits}

Uniform convergence of the base potentials does not control the minimizing balanced metrics themselves: by the cofactor formula \eqref{eq:warped-balanced-cofactor}, their $(n-1)$st powers contain the square of the reciprocal fiber coefficient. Depending on the ratio at the end, the traces of these powers stay bounded or blow up, and in the second case the trace-normalized powers concentrate on fibers over the extreme set.

Keep the sequence $\theta_j\to\theta_0$ of \cref{prop:warped-endpoint-potential-limits}, and let $\phi_0$ be the uniform limit of its base potentials $\phi_j$. Put $A_j=b_j\kappa+\ddc\phi_j$, and write $G_j=\gamma_j^m$, where $\gamma_j$ is the balanced minimizing metric of \cref{thm:warped-gauduchon-volume} for $\theta_j$. Traces are taken with respect to $h=\kappa+w$.

\Needspace{14\baselineskip}
\begin{proposition}\label{prop:warped-minimizer-trace}\index{Gauduchon metric!finite trace boundary}
Let $r>1$ be the ratio $r_0$ or $r_M$ of the end. After discarding finitely many terms, write $L_j=s_j(f+\epsilon_j)$ with $s_j>0$, $s_j\to s_0>0$ and $\epsilon_j\to0^+$. Here $f=a$, $s_j=c_j$ at the first face, and $f=M-a$, $s_j=-c_j$ at the second. Set $J_j=\int_B\sigma/L_j^2$. Then
\begin{equation}\label{eq:warped-minimizer-trace-mass}
 G_j=\frac{\sigma}{nL_j^2}+mA_j^{m-1}\wedge w,
 \qquad t_j:=\int_Xh\wedge G_j=\frac{J_j}{n}+mb_j^{m-1}.
\end{equation}
For $1<r<2$, $t_j\asymp\epsilon_j^{r-2}$; for $r=2$, $t_j\asymp\log(1/\epsilon_j)$. For $r>2$, the entire sequence converges weakly to the finite positive closed current
\[
 G_0=\frac{\sigma}{nL_0^2}+mA_0^{m-1}\wedge w,
 \qquad A_0=b\kappa+\ddc\phi_0,
\]
and $t_j\to t_0=\int_Xh\wedge G_0>0$.
\end{proposition}
\begin{proof}
On $\mathscr S$, $H(c_j,e_j)=nb_j^m$, so $A_j^m=\sigma/(nL_j)$. Substituting this in \eqref{eq:warped-balanced-cofactor} gives the formula for $G_j$. Stokes' theorem on $B$, for the closed form $\kappa$, gives $\int_B\kappa\wedge A_j^{m-1}=b_j^{m-1}$, which proves the trace identity. The factorization of $L_j$ uses $\epsilon_j=e_j/c_j$ at the first end and $\epsilon_j=e_j/(-c_j)-M$ at the second, and the coefficient limits give $s_j\to s_0>0$.

On each tube around an extreme component, $f\asymp\operatorname{dist}(Z_i)^{k_i}$, so the integral of $f^{-2}$ is comparable to
\[
 \int_0^{\varrho_0}\frac{\varrho^{\operatorname{codim}_{\mathbb R}Z_i-1}}{(\varrho^{k_i}+\epsilon_j)^2}\,\mathrm{d}\varrho,
\]
where $\varrho$ is normal radius, $\varrho_0>0$ is small and $r_i=\operatorname{codim}_{\mathbb R}Z_i/k_i$. Splitting the integral at $\varrho=\epsilon_j^{1/k_i}$ gives order $\epsilon_j^{r_i-2}$ for $r_i<2$, order $\log(1/\epsilon_j)$ for $r_i=2$, and a finite limit for $r_i>2$. The finite sum over components is governed by $r=\min_i r_i$. Since $b_j$ and $s_j$ have positive finite limits, these estimates prove the three alternatives for $t_j$.

If $r>2$, then $f^{-2}$ is integrable and dominates $L_j^{-2}$ up to a fixed constant factor, so $\sigma/L_j^2\to\sigma/L_0^2$ in total variation. Since $w\wedge w=0$, we have $A_j^{m-1}\wedge w=(\theta_j+\ddc u_j)^{m-1}\wedge w$, which converges weakly by \cref{prop:warped-endpoint-potential-limits}. This proves $G_j\rightharpoonup G_0$, and $G_0$ is positive and closed as a weak limit of positive closed forms. Pairing with $h$ gives the limit of the traces, which is positive because of the term $\sigma/(nL_0^2)$.
\end{proof}

\Needspace{15\baselineskip}
\begin{proposition}\label{prop:warped-normalized-test-limits}\index{Gauduchon metric!concentrating trace}\index{Gauduchon current!warped fiber limit}
Suppose $1<r\leq2$, and put $T_j=G_j/t_j$. Every weak cluster limit of $(T_j)$ has the form $p_B^*\nu$, where $\nu$ is a probability measure on $B$ supported on the extreme components with ratio $r$. Such a limit is positive and closed, has trace one relative to $h$, and satisfies $\theta_0\wedge p_B^*\nu=0$. If those components reduce to a single point $p$, the whole sequence converges to the current of integration $[\{p\}\times C]$.

The trace-normalized smooth metrics $\eta_j=t_j^{-1/m}\gamma_j$ satisfy
\[
 \int_X\theta_j\wedge\eta_j^m=D_\mu(\eta_j)
                  =\frac{nb_j^m}{t_j}\longrightarrow0.
\]
Both quantities have order $\epsilon_j^{2-r}$ when $1<r<2$, and order $1/\log(1/\epsilon_j)$ when $r=2$. Their quotient remains one.
\end{proposition}
\begin{proof}
Let $\nu_j=\sigma/(J_jL_j^2)$. These are probability measures on $B$, and
\[
 T_j=\frac{J_j}{nt_j}p_B^*\nu_j
                   +\frac m{t_j}A_j^{m-1}\wedge w.
\]
The first coefficient tends to one. The second current is positive with trace $mb_j^{m-1}/t_j\to0$, hence converges weakly to zero. Compactness of probability measures on $B$ and weak continuity of pullback under the submersion $p_B$ therefore show that every cluster limit is $p_B^*\nu$ for a cluster limit $\nu$ of $(\nu_j)$. Away from small tubes around the extreme set, the density $1/L_j^2$ is uniformly bounded, while $J_j\to\infty$; that region receives no mass in the limit. The tube integrals for components with ratio greater than $r$ are $o(J_j)$ by the preceding estimates, so only components with ratio $r$ keep mass.

For a limit $T=p_B^*\nu$, $\kappa\wedge T=0$ for degree reasons on $B$, while $L_0w\wedge T=0$ because $L_0$ vanishes on the support of $\nu$. Its trace is one because $\int_Cw=1$. A top-degree measure on $B$ is closed, so its pullback is closed as well. If the support is the single point $p$, every cluster limit of $(\nu_j)$ is $\delta_p$; compactness then gives convergence of the whole sequence, and $p_B^*\delta_p$ is the current of integration over the fiber $\{p\}\times C$.

By Stokes' theorem and \eqref{eq:warped-minimizer-trace-mass}, $\int_X\theta_j\wedge G_j=\int_X(A_j+L_jw)\wedge G_j=n\int_BA_j^m=nb_j^m$. Since $\gamma_j$ minimizes and $\lambda_\mu(\theta_j)=1$ on $\mathscr S$, also $D_\mu(\gamma_j)=nb_j^m$. Both are multiplied by $t_j^{-1}$ under the trace normalization, and the trace estimates give their limits and rates.
\end{proof}

In the projective strip of \cref{ex:warped-projective-strip}, both ratios are two and each extreme set is a single point. At each of the two ends, the normalized minimizing tests therefore converge to the current of integration over the corresponding fiber $\{p\}\times\mathbb P^1$, and the numerator and the denominator of the quotient both decay like $1/\log(1/\epsilon_j)$. This is convergence of currents; it does not give a pointwise limit of the metrics.
